% ============================================================================
%  CGM Rental — Manual del Panel de Administracion
%  Guia practica para el equipo de Marketing
%  Desarrollado por: Omar Antonio Ccencho Barazorda
%  Cliente: CGM Rental
% ============================================================================
\documentclass[11pt,a4paper]{article}

% ── Paquetes ────────────────────────────────────────────────────────────────
\usepackage[utf8]{inputenc}
\usepackage[T1]{fontenc}
\usepackage[spanish,es-tabla]{babel}
\usepackage[a4paper,margin=2.4cm]{geometry}
\usepackage{xcolor}
\usepackage{graphicx}
\usepackage{titlesec}
\usepackage{enumitem}
\usepackage{tabularx}
\usepackage{longtable}
\usepackage{booktabs}
\usepackage{array}
\usepackage{fancyhdr}
\usepackage{tcolorbox}
\usepackage[hidelinks,colorlinks=true,linkcolor=cgmgreen,urlcolor=cgmgreen]{hyperref}

% ── Colores corporativos CGM ─────────────────────────────────────────────────
\definecolor{cgmgreen}{HTML}{004D3D}
\definecolor{cgmlime}{HTML}{C5E86C}
\definecolor{cgmgray}{HTML}{4A4A4A}

% ── Estilos de titulos ──────────────────────────────────────────────────────
\titleformat{\section}{\color{cgmgreen}\Large\bfseries}{\thesection}{1em}{}
\titleformat{\subsection}{\color{cgmgreen}\large\bfseries}{\thesubsection}{1em}{}
\titleformat{\subsubsection}{\color{cgmgray}\normalsize\bfseries}{\thesubsubsection}{1em}{}

% ── Encabezado y pie de pagina ──────────────────────────────────────────────
\pagestyle{fancy}
\fancyhf{}
\renewcommand{\headrulewidth}{0.4pt}
\renewcommand{\headrule}{\hbox to\headwidth{\color{cgmgreen}\leaders\hrule height \headrulewidth\hfill}}
\fancyhead[L]{\textcolor{cgmgreen}{\textbf{CGM Rental}}}
\fancyhead[R]{\textcolor{cgmgray}{Manual del Panel}}
\fancyfoot[C]{\thepage}

% ── Cajas tematicas ─────────────────────────────────────────────────────────
\newtcolorbox{infobox}[1][]{
    colback=cgmlime!15,
    colframe=cgmgreen,
    boxrule=0.6pt, arc=2pt,
    left=8pt,right=8pt,top=6pt,bottom=6pt, #1
}
\newtcolorbox{warningbox}[1][]{
    colback=red!5,
    colframe=red!50,
    boxrule=0.6pt, arc=2pt,
    left=8pt,right=8pt,top=6pt,bottom=6pt, #1
}
\newtcolorbox{tipbox}[1][]{
    colback=blue!5,
    colframe=blue!40,
    boxrule=0.6pt, arc=2pt,
    left=8pt,right=8pt,top=6pt,bottom=6pt, #1
}
\newtcolorbox{stepbox}[1][]{
    colback=cgmgreen!5,
    colframe=cgmgreen!40,
    boxrule=0.6pt, arc=2pt,
    left=8pt,right=8pt,top=6pt,bottom=6pt, #1
}

% ── Metadatos ───────────────────────────────────────────────────────────────
\title{
    \vspace{-2cm}
    \begin{flushleft}
    \textcolor{cgmgreen}{\rule{\linewidth}{2pt}}\\[0.4cm]
    \textcolor{cgmgreen}{\Huge\bfseries Manual del Panel}\\[0.15cm]
    \textcolor{cgmgreen}{\Huge\bfseries de Administracion}\\[0.25cm]
    \textcolor{cgmgray}{\LARGE Guia para el equipo de Marketing}\\[0.2cm]
    \textcolor{cgmgreen}{\Large CGM Rental}\\[0.2cm]
    \textcolor{cgmlime!70!black}{\rule{\linewidth}{1pt}}
    \end{flushleft}
}

\author{
    \textbf{Desarrollado por:} Omar Antonio Ccencho Barazorda\\
    \textbf{Cliente:} CGM Rental
}

\date{\today}

% ============================================================================
\begin{document}
\maketitle
\thispagestyle{empty}

\begin{infobox}
\textbf{Para quien es esta guia.} Este manual esta escrito para el equipo
de Marketing de CGM Rental. Aqui aprenderas a entrar al panel, gestionar
productos, cambiar banners, publicar noticias del blog y actualizar la
informacion de contacto del sitio web --- todo sin necesidad de saber
programacion.
\end{infobox}

\vspace{0.4cm}
\tableofcontents
\newpage

% ============================================================================
\section{?` Que es el panel de administracion?}
% ============================================================================

El panel de administracion es la \textbf{zona privada} de la pagina web
\texttt{cgmrental.com}. Desde alli puedes:

\begin{itemize}[leftmargin=*]
    \item Ver \textbf{estadisticas} de la pagina (cuantos equipos hay,
        cuantos posts del blog, sucursales, etc.).
    \item \textbf{Agregar, editar o quitar equipos} del catalogo (PE y AR).
    \item \textbf{Cambiar las fotos} de cada equipo.
    \item \textbf{Cambiar los banners} de la pagina principal y de las
        secciones (hero, tarjetas, banners de categoria).
    \item Actualizar el \textbf{telefono, email, redes sociales y direccion}
        que aparecen en el sitio.
    \item Administrar las \textbf{sucursales} (Lima, Piura, Cusco, etc.)
        que se muestran en el mapa.
    \item Publicar y editar \textbf{noticias del blog}.
    \item Ver el \textbf{historial de cambios} (quien hizo que y cuando).
\end{itemize}

\subsection{?` Quien puede entrar?}

Solo las personas autorizadas por CGM Rental. La autorizacion se controla
con tu correo corporativo \texttt{@cgmrental.com}. Si tu correo no esta en
la lista de autorizados, el sistema te dejara afuera aunque sepas la
contrasena de Microsoft.

\begin{tipbox}
\textbf{?` Como me agregan a la lista?} Habla con el responsable tecnico
del sitio o con el administrador del servidor. Solo necesitas darle tu
correo corporativo y te agregan en cuestion de minutos.
\end{tipbox}

% ============================================================================
\section{Como entrar al panel}
% ============================================================================

\subsection{Paso a paso}

\begin{stepbox}
\begin{enumerate}[leftmargin=*]
    \item Abre el navegador (preferiblemente Chrome, Edge o Firefox).
    \item Escribe en la barra de direcciones:\\
        \texttt{https://cgmrental.com/admin/}
    \item Veras la pantalla de bienvenida con un boton \textbf{``Iniciar
        sesion con Microsoft''}.
    \item Haz clic en ese boton.
    \item Microsoft te pedira tu correo corporativo (\texttt{@cgmrental.com})
        y tu contrasena habitual de la empresa.
    \item Si todo esta bien, entraras directamente al \textbf{Dashboard}.
\end{enumerate}
\end{stepbox}

\subsection{No tienes que recordar otra contrasena}

El panel usa la \textbf{misma cuenta de Microsoft} que ya usas para tu
correo de la empresa. No hay usuario y contrasena distintos para el panel.

\subsection{?` Y si me equivoco?}

\begin{itemize}[leftmargin=*]
    \item Si tu correo no esta autorizado, veras un mensaje rojo:
        \textit{``Tu correo no tiene permisos para acceder al panel de
        administracion''}. Solucion: pide al responsable tecnico que te
        agregue a la lista.
    \item Si pones mal la contrasena de Microsoft, esa pantalla la maneja
        Microsoft (no nuestro panel). Recuperala desde tu cuenta corporativa.
\end{itemize}

% ============================================================================
\section{El Dashboard (pantalla de inicio)}
% ============================================================================

Cuando entras al panel, lo primero que ves es el \textbf{Dashboard}. Es
como un tablero con numeros y graficos para que sepas de un vistazo como
esta la pagina.

\subsection{?` Que numeros muestra?}

\begin{itemize}[leftmargin=*]
    \item \textbf{Total de equipos} en el catalogo.
    \item Cuantos estan \textbf{activos} (visibles para el publico) y
        cuantos \textbf{inactivos} (escondidos).
    \item Cuantos tienen \textbf{foto} y cuantos no.
    \item Total de \textbf{noticias} del blog y cuantas estan publicadas.
    \item Cuantas \textbf{sucursales} hay en Peru y en Argentina.
\end{itemize}

\subsection{Graficos}

\begin{itemize}[leftmargin=*]
    \item Equipos por sector (Construccion, Mineria, Agricola, Energia).
    \item Equipos por pais (PE vs AR).
    \item Equipos por tipo de operacion (Alquiler, Venta, Ambos).
    \item Noticias del blog por categoria y por mes.
\end{itemize}

\subsection{Filtro por pais}

En la parte superior puedes elegir ver los datos de:
\textbf{Todos / Peru / Argentina}. Los graficos se actualizan automaticamente
al cambiar la opcion.

\subsection{Actividad reciente}

Al final ves una lista de las ultimas 12 acciones que hicieron los
administradores (subir una foto, editar un equipo, publicar un post, etc.).

% ============================================================================
\section{Productos (Equipos del catalogo)}
% ============================================================================

Esta es la seccion donde gestionas los equipos que se muestran en la
pagina web.

\subsection{Como llegar}
En el menu lateral, haz clic en \textbf{``Productos''}.

\subsection{La pantalla de listado}

Veras una tabla con todos los equipos. Cada fila muestra:
\begin{itemize}[leftmargin=*]
    \item La foto principal del equipo.
    \item Nombre y marca.
    \item Sector y tipo.
    \item Si esta activo o inactivo (con un switch).
    \item En que pais aparece (PE / AR).
    \item Botones para \textbf{Editar} y \textbf{Eliminar}.
\end{itemize}

\subsection{Filtros disponibles}

Arriba de la tabla puedes filtrar por:
\begin{itemize}[leftmargin=*]
    \item \textbf{Buscador:} escribe el nombre, marca o slug del equipo.
    \item \textbf{Pais:} solo PE, solo AR o ambos.
    \item \textbf{Sector:} Construccion, Mineria, Agricola, Energia.
    \item \textbf{Estado:} solo activos o solo inactivos.
\end{itemize}

\subsection{Agregar un equipo nuevo}

\begin{stepbox}
\begin{enumerate}[leftmargin=*]
    \item En la pantalla de Productos, haz clic en \textbf{``Nuevo producto''}.
    \item Llena el formulario con los datos del equipo:
    \begin{itemize}
        \item \textbf{Nombre} (obligatorio): ej.\ ``Excavadora John Deere 210G''.
        \item \textbf{Slug}: dejalo vacio, el sistema lo crea automatico
            a partir del nombre.
        \item \textbf{Marca}: ej.\ ``John Deere''.
        \item \textbf{Tipo}: el tipo de maquina (Excavadora, Cargador
            Frontal, etc.).
        \item \textbf{Sector / Unidad}: a que sector pertenece. Puedes
            poner varios separados por barra vertical:
            \texttt{Construccion|Mineria}.
        \item \textbf{Descripcion}: las caracteristicas del equipo,
            \textbf{separadas por barra vertical}. Ejemplo:\\
            \texttt{Motor 410 HP | Cabina ergonomica | Tanque 200L}
        \item \textbf{Tags}: marca si es para \texttt{alquiler}, para
            \texttt{venta}, o ambos.
        \item \textbf{Activo}: si quieres que se vea en la web ya, marcalo.
        \item \textbf{Mostrar en Argentina}: marcalo si el equipo tambien
            se ofrece en AR.
        \item \textbf{A solicitud}: si el equipo solo se trae bajo pedido,
            marcalo.
    \end{itemize}
    \item Haz clic en \textbf{Guardar}.
    \item Una vez creado, vuelve a entrar a editarlo para subir las
        \textbf{fotos del equipo}.
\end{enumerate}
\end{stepbox}

\subsection{Subir fotos de un equipo}

\begin{stepbox}
\begin{enumerate}[leftmargin=*]
    \item Entra a editar el equipo (boton de lapiz en la lista).
    \item Baja hasta la seccion de \textbf{Imagenes}.
    \item Haz clic en \textbf{``Subir imagen''} y elige el archivo de tu
        computadora.
    \item El sistema automaticamente:
    \begin{itemize}
        \item Convierte la imagen a formato WebP (mas liviana, carga mas
            rapido).
        \item La numera (\texttt{1.webp}, \texttt{2.webp}, etc.). La
            \textbf{primera} imagen sera la \textbf{portada} que se vera
            en el catalogo.
    \end{itemize}
    \item Repite para todas las fotos que quieras subir.
    \item Para borrar una foto, haz clic en el icono de papelera al lado
        de cada imagen.
\end{enumerate}
\end{stepbox}

\begin{tipbox}
\textbf{Tamano maximo:} 10 MB por imagen.\\
\textbf{Formatos aceptados:} WebP, JPG, JPEG, PNG.\\
\textbf{Recomendacion:} fotos en proporcion cuadrada o ligeramente
horizontal, con buena iluminacion y fondo neutro.
\end{tipbox}

\subsection{Activar o desactivar un equipo}

En la lista de productos, cada fila tiene un \textbf{switch}. Haz clic
para encender (activo, visible en la web) o apagar (oculto). El cambio es
inmediato.

\begin{tipbox}
\textbf{Recomendacion.} En vez de \textit{eliminar} un equipo, mejor
\textit{desactivarlo}. Asi puedes volver a activarlo despues sin perder
las fotos ni los datos. Eliminar es definitivo.
\end{tipbox}

\subsection{Eliminar un equipo}

\begin{warningbox}
\textbf{Cuidado.} Eliminar un equipo borra sus datos de la base, pero las
fotos se mantienen en el servidor. Si necesitas recuperarlo despues, hay
que volver a crearlo y subir las fotos otra vez. \textbf{Considera
desactivar en su lugar.}
\end{warningbox}

\subsection{Reconstruir products.json}

En la parte superior de Productos hay un boton:
\textbf{``Reconstruir products.json''}. Esto sirve para casos especiales
cuando algun cambio no se ve correctamente en la web. Si tienes dudas
sobre cuando usarlo, consulta primero al responsable tecnico.

% ============================================================================
\section{Banners}
% ============================================================================

Los banners son las \textbf{imagenes grandes} que se ven en la pagina
principal y en las secciones (la imagen del slider de la home, las
tarjetas, las cabeceras de las categorias, etc.).

\subsection{Como llegar}
Menu lateral $\rightarrow$ \textbf{``Banners''}.

\subsection{Organizacion}

Los banners estan agrupados por seccion para que sea facil encontrarlos:

\begin{itemize}[leftmargin=*]
    \item \textbf{Hero:} la imagen grande de la home (slider).
    \item \textbf{Home Cards:} las tarjetas grandes del home (Construccion,
        Mineria, etc.).
    \item \textbf{Categorias:} las cabeceras de las paginas de categoria.
    \item \textbf{Otros:} contacto, carrito, portal de proveedores, etc.
\end{itemize}

\subsection{Cambiar la imagen de un banner}

\begin{stepbox}
\begin{enumerate}[leftmargin=*]
    \item Encuentra el banner que quieres cambiar en la lista.
    \item Haz clic en \textbf{``Reemplazar imagen''}.
    \item Elige el archivo nuevo desde tu computadora.
    \item Listo. El cambio se ve inmediatamente en la web.
\end{enumerate}
\end{stepbox}

\subsection{Tamanos recomendados}

Cada tipo de banner tiene un tamano ideal para que se vea bien sin
recortarse:

\begin{center}
\begin{tabular}{l l l}
\toprule
\textbf{Tipo de banner} & \textbf{Tamano recomendado} & \textbf{Proporcion} \\
\midrule
Hero (slider home)    & 2700 $\times$ 1000 px  & panoramica \\
Categoria (cabecera)  & 1600 $\times$ 600 px   & panoramica \\
Home Card (estandar)  & 800 $\times$ 700 px    & casi cuadrada \\
Home Card leasing     & 1400 $\times$ 580 px   & rectangular \\
Contacto (banner)     & 2700 $\times$ 1000 px  & panoramica \\
Contacto lateral      & 1100 $\times$ 800 px   & rectangular \\
Carrito               & 1400 $\times$ 580 px   & rectangular \\
Portal proveedores    & 1100 $\times$ 800 px   & rectangular \\
\bottomrule
\end{tabular}
\end{center}

\subsection{Banner distinto para Peru y Argentina}

A veces quieres que el banner se vea distinto en cada pais (por ejemplo,
un banner promocional especifico para Argentina). Para eso:

\begin{stepbox}
\begin{enumerate}[leftmargin=*]
    \item En el banner que quieras, busca el boton \textbf{``Subir version
        pais''}.
    \item Elige Peru o Argentina y sube el archivo.
    \item Listo: ese pais vera tu version especifica; el otro seguira
        viendo la version general.
\end{enumerate}
\end{stepbox}

Para quitar la version pais y volver a usar la general, hay un boton
\textbf{``Eliminar version pais''}.

\subsection{Por que no puedo borrar un banner}

\begin{warningbox}
Los banners principales solo se pueden \textbf{reemplazar}, no eliminar
del todo. Esto evita que se rompa la pagina si por error se borra una
imagen que el sitio espera encontrar. Solo se pueden eliminar los
\textit{archivos huerfanos} (imagenes viejas que ya no estan en uso).
\end{warningbox}

% ============================================================================
\section{Contacto y Redes Sociales}
% ============================================================================

Aqui se cambia la \textbf{informacion de contacto} que aparece en el
footer y en la pagina de Contacto del sitio.

\subsection{Como llegar}
Menu lateral $\rightarrow$ \textbf{``Contacto''}.

\subsection{Que puedes editar}

Tienes una pestana para \textbf{Peru} y otra para \textbf{Argentina}. En
cada una puedes cambiar:

\begin{itemize}[leftmargin=*]
    \item Telefono visible.
    \item Link de WhatsApp.
    \item Email.
    \item Facebook, Instagram, LinkedIn, YouTube.
    \item Direccion principal.
    \item Ciudad principal.
    \item Video destacado de YouTube.
\end{itemize}

\subsection{Como cambiar el telefono o WhatsApp}

\begin{stepbox}
\begin{enumerate}[leftmargin=*]
    \item Entra a la pestana del pais (Peru o Argentina).
    \item Cambia el valor en el campo correspondiente.
    \item Haz clic en \textbf{``Guardar''} al final.
\end{enumerate}
\end{stepbox}

El cambio se ve en cuanto recargas la pagina publica.

\subsection{Video de YouTube destacado}

Si quieres cambiar el video que se muestra en el home:

\begin{stepbox}
\begin{enumerate}[leftmargin=*]
    \item Copia el link del video de YouTube (cualquier formato funciona:
        \texttt{youtube.com/watch?v=XXX}, \texttt{youtu.be/XXX}, etc.).
    \item Pegalo en el campo \textbf{``Video YouTube''}.
    \item Guarda. El sistema extrae automaticamente el codigo del video y
        lo muestra en el home.
\end{enumerate}
\end{stepbox}

% ============================================================================
\section{Sucursales}
% ============================================================================

Aqui gestionas las sucursales de CGM Rental que aparecen en el mapa de
\textbf{``Nuestros Locales''}.

\subsection{Como llegar}
Menu lateral $\rightarrow$ \textbf{``Sucursales''}.

\subsection{Tipos de sucursal}

\begin{itemize}[leftmargin=*]
    \item \textbf{Principal:} sede principal del pais (1 por pais).
    \item \textbf{Sucursal:} sucursal con operacion completa.
    \item \textbf{Atencion:} punto de atencion sin operacion logistica.
\end{itemize}

\subsection{Crear una sucursal nueva}

\begin{stepbox}
\begin{enumerate}[leftmargin=*]
    \item Haz clic en \textbf{``Nueva sucursal''}.
    \item Elige el pais (Peru o Argentina).
    \item Llena los datos:
    \begin{itemize}
        \item Nombre (ej.\ ``Piura'').
        \item Tipo (principal / sucursal / atencion).
        \item Direccion completa.
        \item Latitud y Longitud (busca en Google Maps el lugar, haz
            clic derecho y copia las coordenadas).
        \item Link de Google Maps (opcional pero recomendado, asi el
            boton ``Como llegar'' funciona).
        \item Telefono (opcional).
        \item Orden (numero para ordenar en la lista; mas bajo aparece
            primero).
    \end{itemize}
    \item Guarda.
\end{enumerate}
\end{stepbox}

\subsection{Editar o eliminar}

Cada fila tiene botones de editar y eliminar. Al editar, cambias los
mismos datos del paso anterior.

\begin{tipbox}
\textbf{?` Como saco las coordenadas de Google Maps?}\\
1. Abre Google Maps y busca el lugar.\\
2. Haz clic derecho en el punto exacto sobre el mapa.\\
3. Veras dos numeros (ej.\ \texttt{-12.2773, -76.8864}). El primero es
\textbf{latitud}, el segundo \textbf{longitud}.
\end{tipbox}

% ============================================================================
\section{Blog / Novedades}
% ============================================================================

Esta es la seccion para publicar \textbf{noticias, articulos y entrevistas}
en \texttt{/novedades/}.

\subsection{Como llegar}
Menu lateral $\rightarrow$ \textbf{``Blog''}.

\subsection{Crear un post nuevo}

\begin{stepbox}
\begin{enumerate}[leftmargin=*]
    \item Haz clic en \textbf{``Nuevo post''}.
    \item Llena los campos:
    \begin{itemize}
        \item \textbf{Titulo} (obligatorio).
        \item \textbf{Slug}: dejalo vacio, se genera automatico desde el
            titulo.
        \item \textbf{Categoria}: \texttt{articulos}, \texttt{entrevistas},
            etc.
        \item \textbf{Fecha} en formato \texttt{YYYY-MM-DD} (ej.\
            \texttt{2026-05-15}).
        \item \textbf{Extracto}: el resumen breve que aparece en la lista
            de noticias (1--2 lineas).
        \item \textbf{Contenido}: el cuerpo del articulo en HTML. Para
            titulos usa \texttt{<h3>...</h3>} y para parrafos
            \texttt{<p>...</p>}.
        \item \textbf{Video URL}: si tienes un video de YouTube relacionado,
            pega el link.
        \item \textbf{Imagen}: sube la imagen principal del post.
        \item \textbf{Activo}: marcalo para que se publique.
        \item \textbf{Mostrar en PE} y \textbf{Mostrar en AR}: elige en
            que paises se vera. Puedes marcar uno o ambos.
    \end{itemize}
    \item Guarda.
\end{enumerate}
\end{stepbox}

\subsection{Activar / desactivar un post}

Al lado de cada post hay un switch. Igual que con los productos:
encendido = visible en la web, apagado = oculto.

\subsection{Eliminar un post}

Boton de papelera al final de cada fila. Igual que con productos, eliminar
es definitivo.

\begin{warningbox}
\textbf{Importante.} Si guardas un post \textit{activo} pero \textbf{no}
marcas ningun pais (PE o AR), el sistema te avisara pero igual lo guardara.
El post no aparecera en la web hasta que actives al menos un pais.
\end{warningbox}

% ============================================================================
\section{Historial de cambios}
% ============================================================================

Aqui ves el \textbf{registro de todas las acciones} que se han hecho en
el panel: quien hizo que y cuando.

\subsection{Como llegar}
Menu lateral $\rightarrow$ \textbf{``Historial''}.

\subsection{Para que sirve}

\begin{itemize}[leftmargin=*]
    \item Si algo desaparecio de la web y no sabes por que, aqui ves quien
        lo cambio o lo desactivo.
    \item Si quieres auditar quien esta haciendo cambios.
    \item Si quieres revisar tu propia actividad.
\end{itemize}

\subsection{Filtrar por usuario}

Puedes filtrar por el correo de un usuario especifico para ver solo sus
acciones.

% ============================================================================
\section{Tareas comunes paso a paso}
% ============================================================================

Aqui tienes las tareas mas frecuentes resueltas de principio a fin.

\subsection{?` Como cambiar la foto del slider de la home?}

\begin{stepbox}
\begin{enumerate}[leftmargin=*]
    \item Entra al panel.
    \item Menu $\rightarrow$ Banners.
    \item Busca el grupo \textbf{``Hero''}.
    \item Encuentra el slide que quieres cambiar.
    \item Clic en \textbf{Reemplazar imagen}.
    \item Sube tu nueva imagen (preferiblemente 2700 $\times$ 1000 px).
    \item Listo, recarga la home y la veras.
\end{enumerate}
\end{stepbox}

\subsection{?` Como agregar un equipo que solo va a Argentina?}

\begin{stepbox}
\begin{enumerate}[leftmargin=*]
    \item Menu $\rightarrow$ Productos $\rightarrow$ Nuevo producto.
    \item Llena nombre, marca, tipo, sector, descripcion.
    \item Marca el checkbox \textbf{``Mostrar en Argentina''}.
    \item \textbf{Desmarca} \textbf{``Mostrar en Peru''}.
    \item Marca \textbf{``Activo''}.
    \item Guarda y sube las fotos.
\end{enumerate}
\end{stepbox}

\subsection{?` Como publicar una noticia en ambos paises?}

\begin{stepbox}
\begin{enumerate}[leftmargin=*]
    \item Menu $\rightarrow$ Blog $\rightarrow$ Nuevo post.
    \item Llena titulo, fecha, categoria, extracto y contenido.
    \item Sube la imagen principal.
    \item Marca \textbf{Activo}, \textbf{Mostrar en PE} y \textbf{Mostrar
        en AR}.
    \item Guarda.
\end{enumerate}
\end{stepbox}

\subsection{?` Como cambiar el WhatsApp de Argentina?}

\begin{stepbox}
\begin{enumerate}[leftmargin=*]
    \item Menu $\rightarrow$ Contacto.
    \item Cambia a la pestana \textbf{Argentina}.
    \item Edita el campo \textbf{``Link de WhatsApp''}.
    \item Guarda.
\end{enumerate}
\end{stepbox}

\subsection{?` Como ocultar un equipo temporalmente?}

\begin{stepbox}
\begin{enumerate}[leftmargin=*]
    \item Menu $\rightarrow$ Productos.
    \item Encuentra el equipo en la lista.
    \item Apaga el switch \textbf{Activo}.
    \item Listo. El equipo desaparece de la web pero no se borra.
\end{enumerate}
\end{stepbox}

\subsection{?` Como abrir una sucursal nueva?}

\begin{stepbox}
\begin{enumerate}[leftmargin=*]
    \item Busca las coordenadas en Google Maps (clic derecho sobre el
        local).
    \item Menu $\rightarrow$ Sucursales $\rightarrow$ Nueva sucursal.
    \item Llena nombre, tipo (``sucursal'' generalmente), direccion, lat
        y lng.
    \item Pega el link de Google Maps en \textbf{``maps\_url''}.
    \item Guarda.
\end{enumerate}
\end{stepbox}

% ============================================================================
\section{Problemas comunes y como resolverlos}
% ============================================================================

\subsection{``Sube la imagen pero no la veo en la web''}

Lo mas probable es que el navegador este mostrando la version antigua que
tiene guardada en cache.

\begin{tipbox}
\textbf{Solucion:} en la pagina publica presiona \textbf{Ctrl + F5} (en
Windows) o \textbf{Cmd + Shift + R} (en Mac). Esto fuerza al navegador a
descargar la version mas reciente.
\end{tipbox}

\subsection{``Mi correo no tiene permisos''}

Tu correo no esta en la lista de autorizados.

\textbf{Solucion:} contacta al responsable tecnico para que te agregue.

\subsection{``Edite un equipo y al rato volvio al estado anterior''}

Es un caso raro de desincronizacion entre la base de datos y un archivo
de respaldo.

\textbf{Solucion:} ve a Productos y haz clic en \textbf{``Reconstruir
products.json''}. Si el problema persiste, contacta al responsable
tecnico.

\subsection{``El video de YouTube no aparece en el home''}

Posiblemente el campo se guardo con el link mal copiado.

\textbf{Solucion:} ve a Contacto, copia y pega el link completo del
video de YouTube en \textbf{``Video YouTube''} y guarda. El sistema lo
procesa automaticamente.

\subsection{``Quiero borrar un banner pero no puedo''}

Es a proposito. Los banners principales solo se reemplazan, nunca se
borran del todo, para que la pagina no se rompa.

\textbf{Solucion:} si quieres ``ocultar'' un banner, sube una imagen
neutra o el banner anterior.

\subsection{``Subi una imagen muy grande y me dio error''}

El limite es 10 MB por imagen.

\textbf{Solucion:} comprime la imagen en algun servicio gratuito como
TinyPNG, Squoosh o Photoshop, y vuelve a subirla.

% ============================================================================
\section{Cosas que el panel todavia no hace}
% ============================================================================

Para que tengas claras las limitaciones actuales:

\begin{itemize}[leftmargin=*]
    \item \textbf{No muestra las cotizaciones recibidas.} Las cotizaciones
        que envia la gente desde la web llegan por correo a Ventas, pero
        el panel no tiene una pantalla para verlas en lista.
    \item \textbf{No muestra los proveedores ni las denuncias.} Igual que
        las cotizaciones: llegan por correo, pero no hay pantalla en el
        panel.
    \item \textbf{No tiene roles diferentes.} Todas las personas
        autorizadas pueden hacer todo. No hay un rol ``solo blog'' o
        ``solo lectura''.
    \item \textbf{No tiene papelera.} Si eliminas un producto o un post,
        no se puede recuperar desde el panel. Mejor desactiva en vez de
        eliminar.
    \item \textbf{No tiene vista previa.} Para ver como queda un producto
        antes de publicarlo, hay que activarlo y abrir la pagina publica.
    \item \textbf{No tiene edicion masiva.} No se pueden marcar varios
        productos a la vez para desactivar o cambiar de sector.
\end{itemize}

% ============================================================================
\section{Buenas practicas}
% ============================================================================

\begin{itemize}[leftmargin=*]
    \item \textbf{Antes de eliminar, desactiva.} Es facil deshacer una
        desactivacion; no se puede deshacer una eliminacion.
    \item \textbf{Usa nombres claros} en los productos y posts. Eso ayuda
        al SEO (Google) y a tus propios usuarios.
    \item \textbf{Sube siempre la mejor foto como primera}, porque sera
        la portada que se ve en el catalogo.
    \item \textbf{Antes de publicar un post, leelo entero}, porque una vez
        publicado los buscadores pueden indexarlo al toque.
    \item \textbf{Manten el historial limpio:} usa tu propio correo, no
        compartas la cuenta con otros. Asi en el historial queda claro
        quien hizo cada cambio.
    \item \textbf{Si tienes dudas, no toques.} Mejor preguntar al
        responsable tecnico que ``probar a ver que pasa'' en produccion.
\end{itemize}

% ============================================================================
\section{Glosario simple}
% ============================================================================

\begin{description}[leftmargin=*]
    \item[\textbf{Activo / Inactivo}] Si esta encendido (visible en la web)
        o apagado (oculto).
    \item[\textbf{Banner}] Una de las imagenes grandes del sitio (slider,
        cabeceras, tarjetas).
    \item[\textbf{Catalogo}] La lista de equipos publicados en
        \texttt{/categoria-producto/}.
    \item[\textbf{Dashboard}] La pantalla de inicio del panel con numeros
        y graficos.
    \item[\textbf{Equipo}] Cada maquina del catalogo (excavadora, rodillo,
        etc.). Tambien llamado ``producto''.
    \item[\textbf{Hero}] La imagen / slider grande de la pagina principal.
    \item[\textbf{Pais}] PE = Peru, AR = Argentina. Cada uno tiene su
        propia vitrina en \texttt{cgmrental.com/pe/} y \texttt{/ar/}.
    \item[\textbf{Slug}] El nombre corto del producto en la URL (ej.\
        \texttt{excavadora-210g}). Se genera automatico.
    \item[\textbf{Slot}] Cada ``espacio'' fijo donde va un banner. Cuando
        cambias la imagen, el slot sigue siendo el mismo.
    \item[\textbf{Tag}] La etiqueta de operacion: \texttt{alquiler} o
        \texttt{venta}.
    \item[\textbf{Unidad / Sector}] El area del negocio: Construccion,
        Mineria, Agricola o Energia.
    \item[\textbf{Webp}] Un formato de imagen mas liviano que JPG/PNG. El
        sistema convierte automaticamente todo a WebP.
\end{description}

\vspace{1cm}
\begin{infobox}
\textbf{?` Tienes dudas o algo no funciona?}\\
Si encuentras un problema, un error o quieres una funcion nueva en el
panel, contacta al desarrollador del proyecto:\\
\textbf{Omar Antonio Ccencho Barazorda}.
\end{infobox}

\end{document}
